% SGA 3, 第 VI_B 篇, 第 1 轮英文重构.
% 范围: 命题 3.11, 第 4 节标题以及命题 4.1.
% 底本: Polo--Gille Exp6B-13oct24.pdf, 分篇第 22 页
% (合订本第 363 页; 原版页边印刷页码 347--348).
% OCR 与 Reinhold 英文 Markdown 仅为比对材料.

\medskip
\phantomsection\label{prop:sga3-VIB-3.11}%
\noindent\textbf{命题 3.11.---}%
\footnote{\label{note:sga3-VIB-26}编者注: 比较
VI\(_A\), \SGARefLink{par:sga3-VIB-5.6}{5.6}.}
\emph{设 \(u:G\to H\) 是局部有限呈示 \(S\)-群之间的态射.
若 \(u\) 平坦, 则诱导映射}
\[
  u^0:G^0\longrightarrow H^0
\]
\emph{满射. 若 \(G\) 在 \(S\) 上平坦, 且 \(H\) 的纤维既约, 则逆命题成立.}

假设 \(u\) 平坦. 则对每个 \(s\in S\), \(u_s\) 平坦且局部有限呈示,
因而开 (EGA~IV\(_2\), 2.4.6). 因此由 \SGARefLink{cor:sga3-VIB-1.3.2}{1.3.2},
\[
  u_s^0:\underline{G_s^0}\longrightarrow\underline{H_s^0}
\]
满射.%
\footnote{\label{note:sga3-VIB-27}编者注: 这里通过引用
\SGARefLink{cor:sga3-VIB-1.3.2}{1.3.2} 缩短了原文.
还省去了对一般平坦性定理 (EGA~IV\(_2\), 6.9.1) 的引用,
简化了下面的论证, 因为此处不需要该定理.}
于是映射
\[
  u^0:G^0\longrightarrow H^0
\]
满射.

反过来, 假设 \(u^0:G^0\to H^0\) 满射,
\(G\) 在 \(S\) 上平坦, 且 \(H\) 的纤维既约.
则对每个 \(s\in S\), 态射
\[
  u_s^0:\underline{G_s^0}\longrightarrow\underline{H_s^0}
\]
满射, 因而在位于 \(\underline{H_s^0}\) 的泛点 \(\eta\)
上的各点处平坦, 因为
\[
  \mathcal O_{\underline{H_s^0},\eta}
\]
是域. 由 \hyperref[prop:sga3-VIB-1.3]{1.3} 可知 \(u_s\) 平坦.
因此由逐纤维平坦性判据 (EGA~IV\(_3\), 11.3.11), \(u\) 平坦.

\section*{4. 局部有限呈示群的纤维维数}
\phantomsection\label{sec:sga3-VIB-4}%
\addcontentsline{toc}{section}{4. 局部有限呈示群的纤维维数}

\medskip
\phantomsection\label{prop:sga3-VIB-4.1}%
\noindent\textbf{命题 4.1.---}
\emph{设 \(G\) 是局部有限型 \(S\)-概形,
带有一个 \(S\)-截面 \(\epsilon\), 并假设对每个 \(s\in S\),}
\[
  \dim G_s=\dim_{\epsilon(s)}G_s
\]
\emph{\textup{(当 \(G\) 是 \(S\)-群时即有此性质;
\hyperref[lem:sga3-VIB-1.5]{1.5})}.}
\begin{enumerate}
\item[\textup{\SGAInlineRefTarget{prop:sga3-VIB-4.1:item:i}(i)}] \emph{函数}
  \[
    s\longmapsto\dim G_s
  \]
  \emph{在 \(S\) 上上半连续.}
\item[\textup{\SGAInlineRefTarget{prop:sga3-VIB-4.1:item:ii}(ii)}]
  \emph{若 \(G\) 在 \(S\) 上还局部有限呈示, 则此函数局部可构造.}
\end{enumerate}

设
\[
  \pi:G\longrightarrow S
\]
是结构态射. Chevalley 半连续性定理 (EGA~IV\(_3\), 13.1.3) 断言函数
\[
  x\longmapsto\dim_x\pi^{-1}\bigl(\pi(x)\bigr)
\]
在 \(G\) 上上半连续. 但对每个 \(s\in S\),
\[
  \dim G_s
    =\dim\pi^{-1}(s)
    =\dim_{\epsilon(s)}
       \pi^{-1}\bigl(\pi(\epsilon(s))\bigr).
\]
由于映射 \(s\mapsto\epsilon(s)\) 在 \(S\) 上连续, 复合函数
\[
  s\longmapsto\dim G_s
\]
在 \(S\) 上上半连续.

现假设 \(G\) 在 \(S\) 上局部有限呈示. 如上论证,
为证明 \(s\mapsto\dim G_s\) 局部可构造, 只须证明函数
\[
  x\longmapsto\dim_x\pi^{-1}\bigl(\pi(x)\bigr)
\]
在 \(G\) 上局部可构造, 而这由 EGA~IV\(_3\), 9.9.1 得到.
